\section{Полиномы Жегалкина (три способа получения). Линейное представление булевых функций.}

\subsection*{Полином Жегалкина}

\paragraph{Определение}
\textbf{Полиномом Жегалкина} называется выражение (полином) над полем $\mathbb{Z}_2$, то есть формула в базисе $\{\oplus, \land, 1\}$. 
В такой формуле отсутствуют отрицания и дизъюнкции, а операции подчиняются законам арифметики по модулю 2:
$$x \oplus x = 0, \quad x \oplus 0 = x, \quad x \land x = x, \quad x \land 1 = x.$$

\paragraph{Теорема о единственности}
Всякая булева функция единственным образом (с точностью до порядка слагаемых) представляется в виде полинома Жегалкина.

\subsection*{Три способа получения полинома}

\begin{enumerate}
    \item \textbf{Метод эквивалентных преобразований:}
    Берется любая формула (например, СДНФ) и в ней производятся замены:
    \begin{itemize}
        \item $\neg x \to x \oplus 1$
        \item $x \lor y \to x \oplus y \oplus (x \land y)$
    \end{itemize}
    Затем раскрываются скобки по дистрибутивности, и слагаемые с четным числом вхождений взаимно уничтожаются ($A \oplus A = 0$).

    \item \textbf{Метод неопределенных коэффициентов:}
    Полином ищется в общем виде: $f(x_1, x_2) = a_0 \oplus a_1 x_1 \oplus a_2 x_2 \oplus a_{12} x_1 x_2$.
    Подставляя все $2^n$ наборов значений переменных, получаем систему линейных уравнений относительно коэффициентов $a_i \in \{0, 1\}$. Система всегда имеет треугольный вид и легко решается.

    \item \textbf{Метод треугольника (сумматорный метод):}
    Строится таблица, в первом столбце которой записываются значения функции из таблицы истинности. Каждый следующий столбец получается путем сложения по модулю 2 соседних элементов предыдущего столбца. Верхние элементы полученных столбцов будут являться коэффициентами полинома Жегалкина при соответствующих конъюнкциях в лексикографическом порядке.
\end{enumerate}

\subsection*{Линейное представление}

Булева функция называется \textbf{линейной}, если её полином Жегалкина не содержит произведений переменных (т.е. его степень $\le 1$):
$$f(x_1, \dots, x_n) = a_0 \oplus a_1 x_1 \oplus \dots \oplus a_n x_n.$$
Линейные функции играют ключевую роль в теории кодирования и криптографии.